\documentclass[11pt,a4paper]{article}
\usepackage[margin=1in]{geometry}
\usepackage[T1]{fontenc}
\usepackage{lmodern}
\usepackage{microtype}
\usepackage{amsmath,amssymb}
\usepackage{xcolor}
\usepackage{hyperref}
\hypersetup{colorlinks=true,linkcolor=blue!50!black,urlcolor=blue!50!black}
\title{A Strongly Typed C-like Subset Lowered to AArch64}
\author{Typed checking, normalization, control flow, and explicit effects}
\date{16 August 2026}
\begin{document}
\maketitle
\begin{abstract}
This note extends the compact Python-like compiler into a conventional strongly typed C-like subset. The supported core contains integers, booleans, strings, fixed-size integer arrays, function definitions with typed parameters, conditionals, explicit array indexing, and one documented Linux AArch64 system call. A static checker validates the indexed result type before normalization and lowering. The emitter uses a small deterministic register convention and stack spills for binary expressions; it does not introduce dynamic language features or hidden runtime semantics.
\end{abstract}
\section{Language and types}
Types are $\mathsf{int}$, $\mathsf{bool}$, $\mathsf{string}$, fixed arrays $\mathsf{int}[n]$, and $\mathsf{void}$. Expressions include literals, variables, arrays, indexing, arithmetic, equality, boolean values, conditionals, and calls. A function has the type
\[
(\tau_1,\ldots,\tau_n)\to\tau.
\]
The checker maintains a value environment and a function environment. It rejects an expression unless every operator, branch, array index, and call has the required type. In particular, both branches of an `if` must have the same result type.

\section{Normalization invariant}
The normalizer folds static arithmetic and boolean expressions before code generation. For example, the condition `2 == 2` becomes `true`, although the generated sample retains the conditional to demonstrate control-flow lowering. A typed transformation preserves the result type:
\[
\Gamma\vdash e:\tau\quad\Longrightarrow\quad
\Gamma\vdash\mathsf{normalize}(e):\tau.
\]
The emitter consumes only checked expressions, so normalization cannot turn an integer expression into a string or an array into a scalar.

\section{ABI subset}
The compiler documents a small AArch64 convention:
\begin{itemize}
\item integer and boolean values use general registers;
\item the first three call arguments use $x_0,x_1,x_2$;
\item function results use $x_0$;
\item $x_{29}$ is the frame pointer and $x_{30}$ the link register;
\item binary-expression temporaries spill through 16-byte stack slots;
\item `sys\_write(fd, buffer, length)` uses Linux syscall number 64 in $x_8$ and executes `svc \#0`.
\end{itemize}
No implicit heap, exception system, closure conversion, or dynamic dispatch is added.

\section{Sample program}
The AST-built sample contains a two-argument integer function `add`, a fixed integer array and index, a string passed to `sys\_write`, a boolean equality condition, and an integer conditional return. The checker reports `typecheck: PASS`. The generated assembly contains `add`, `main`, `cset x0, eq`, conditional branches, `sys\_write`, and `.rodata` entries for the array and string.

\section{Implementation}
The implementation is \texttt{typed\_subset.cpp}:
\begin{verbatim}
g++ -std=c++23 -Wall -Wextra -pedantic -O2 typed_subset.cpp -o typed_subset
./typed_subset > typed_subset.s
\end{verbatim}
The typed AST is explicit and immutable. The code generator uses a small register map for parameters and locals, with stack preservation around binary operations. This is a simple register-allocation policy, not a global graph-coloring allocator; its predictability is useful while the type and lowering invariants are being established.

\section{Boundaries and extensions}
The compiler intentionally excludes implicit coercions, objects, exceptions, garbage collection, variadic calls, and undocumented runtime helpers. The next principled additions are a parser for the already-defined AST, explicit effect types for I/O, ABI-checked function declarations, and a verified lowering relation. Each should preserve
\[
\Gamma\vdash e:\tau
\Longrightarrow
\mathsf{normalize}(e):\tau
\Longrightarrow
\mathsf{emit}(e):\mathsf{Asm}(\tau).
\]
\end{document}
